% !TEX program = xelatex
\documentclass[8pt,aspectratio=169]{beamer}

\usepackage{tatalk}
\addbibresource{references.bib}

\renewcommand{\outlinefontsize}{\small}
\renewcommand{\outlineitemskip}{3pt}

%%%%%%%%%%%%%%%%%%%%%%%%%%%%%%%%%%%%%%%%%%%%%%%%%%%%%%%%%%%%%%%%%%%%%%%%%%%%%
\title{Distant Viewing : Tutoriel 2}
\author{Taylor Arnold}
\email{tarnold2@richmond.edu}
\institute{Distant Viewing Lab}
\date{Modèles avancés}

\begin{document}

%%%%%%%%%%%%%%%%%%%%%%%%%%%%%%%%%%%%%%%%%%%%%%%%%%%%%%%%%%%%%%%%%%%%%%%%%%%%%
\customtitleframe

%%%%%%%%%%%%%%%%%%%%%%%%%%%%%%%%%%%%%%%%%%%%%%%%%%%%%%%%%%%%%%%%%%%%%%%%%%%%%
\framesection{Introduction}
\begin{rightframe}{Bienvenue}
\splitcols{%
  Ces diapositives accompagnent les tutoriels Python que nous avons conçus pour introduire les fondements théoriques et méthodologiques du distant viewing, ainsi qu'un aperçu des outils du Distant Viewing Toolkit.
  
  Elles contiennent des notes, références et schémas supplémentaires que nous trouvons utiles lors des ateliers en présentiel. La plupart des notes techniques et du matériel se trouvent directement dans les notebooks.

  Pour en savoir plus sur notre travail, rendez-vous sur le site du Distant Viewing Lab à l'adresse https://distantviewing.org, ou écrivez-nous directement à tarnold2@richmond.edu.
}{%
  \imgcaption{media/image1.png}{Un télescope surplombant Paris, en clin d'œil à la méthodologie du distant viewing.}
}
\end{rightframe}

%%%%%%%%%%%%%%%%%%%%%%%%%%%%%%%%%%%%%%%%%%%%%%%%%%%%%%%%%%%%%%%%%%%%%%%%%%%%%
\begin{outlineframe}{Plan}
  \rmitem{2.1 Préparation des données}
  \rmdim{2.2 Détection d'objets}
  \rmdim{2.3 Détection à vocabulaire ouvert}
  \rmdim{2.4 Grounding DINO + TrOCR}
  \rmdim{2.5 Segment Anything}
  \rmdim{2.6 SAM + Grounding DINO}
  \rmdim{2.7 Estimation de profondeur}
  \rmdim{2.8 Plongements d'images}
  \rmdim{2.9 Modèles contrastifs}
  \rmdim{2.10 Détection et reconnaissance de visages}
  \rmdim{2.11 Estimation de pose}
  \rmdim{2.12 Modèles vision-langage (VLM)}
  \rmdim{2.13 VLM avec sortie structurée}
  \rmdim{2.14 Conclusion}
\end{outlineframe}

%%%%%%%%%%%%%%%%%%%%%%%%%%%%%%%%%%%%%%%%%%%%%%%%%%%%%%%%%%%%%%%%%%%%%%%%%%%%%
\framesection{2.1 Préparation des données}
\begin{rightframe}{Python et Colab}
\splitcols{%
  Ce tutoriel s'appuie sur l'environnement Google Colab, qui permet d'utiliser Python directement dans le navigateur, sans installer de logiciel sur sa machine. Colab est gratuit pour les petits projets, et des formules payantes donnent accès à davantage de ressources de calcul.

  Nous utiliserons trois modules Python courants pour l'analyse de données : NumPy, Polars et plotnine.
  
  Toutes les bibliothèques utilisées dans le tutoriel sont disponibles sous licence libre. Le code peut aussi être exécuté sur votre propre machine, moyennant une configuration un peu plus poussée.
}{%
  \imgcaption{media/imageD.png}{Logos de Google Colab, Python, Pandas, NumPy et Matplotlib.}
}
\end{rightframe}

%%%%%%%%%%%%%%%%%%%%%%%%%%%%%%%%%%%%%%%%%%%%%%%%%%%%%%%%%%%%%%%%%%%%%%%%%%%%%
\begin{rightframe}{Exécuter le notebook}
\splitcols{%
  Lorsqu'on ouvre le notebook du tutoriel dans Colab, une fenêtre similaire à celle de droite s'affiche dans le navigateur. Le notebook est consultable par tous ; pour exécuter le code, il faut se connecter à un compte Google.

  Remarque : vous pouvez modifier le code et le texte localement dans votre navigateur. Ces modifications ne seront pas enregistrées pour les autres. N'hésitez donc pas à expérimenter avec le code à tout moment !

  Une fois connecté à un compte Google, vous pourrez exécuter les blocs de code (les portions sur fond gris). Pour cela, passez la souris sur le fond gris et cliquez sur le bouton de lecture indiqué par la flèche rouge. Tout au long du tutoriel, il faut exécuter chaque bloc dans l'ordre.
}{%
  \imgcaption{media/image9.png}{Capture d'écran de l'interface Colab exécutant le code de configuration du DVT.}
}
\end{rightframe}

%%%%%%%%%%%%%%%%%%%%%%%%%%%%%%%%%%%%%%%%%%%%%%%%%%%%%%%%%%%%%%%%%%%%%%%%%%%%%
\begin{rightframe}{Présentation du corpus}
\splitcols{%
  Dans ce tutoriel, nous allons explorer une collection d'affiches issues des 100 films les plus rentables de chaque année entre 1970 et 2019. Notre corpus comprend près de 5000 images (nous n'avons pas pu retrouver un petit nombre d'affiches anciennes).

  Notre objectif est de comprendre comment la couleur et la composition des affiches (1) ont évolué dans le temps et (2) sont utilisées pour véhiculer (ou détourner) les conventions de genre. Nous voulons mobiliser des techniques computationnelles pour identifier et comprendre des motifs à travers les milliers d'images de notre collection.
}{%
  \imgcaption{media/image10.png}{Exemples d'affiches de films à travers les décennies et les genres du jeu de données.}
}
\end{rightframe}

%%%%%%%%%%%%%%%%%%%%%%%%%%%%%%%%%%%%%%%%%%%%%%%%%%%%%%%%%%%%%%%%%%%%%%%%%%%%%
\begin{outlineframe}{Plan}
  \rmdim{2.1 Préparation des données}
  \rmitem{2.2 Détection d'objets}
  \rmdim{2.3 Détection à vocabulaire ouvert}
  \rmdim{2.4 Grounding DINO + TrOCR}
  \rmdim{2.5 Segment Anything}
  \rmdim{2.6 SAM + Grounding DINO}
  \rmdim{2.7 Estimation de profondeur}
  \rmdim{2.8 Plongements d'images}
  \rmdim{2.9 Modèles contrastifs}
  \rmdim{2.10 Détection et reconnaissance de visages}
  \rmdim{2.11 Estimation de pose}
  \rmdim{2.12 Modèles vision-langage (VLM)}
  \rmdim{2.13 VLM avec sortie structurée}
  \rmdim{2.14 Conclusion}
\end{outlineframe}

%%%%%%%%%%%%%%%%%%%%%%%%%%%%%%%%%%%%%%%%%%%%%%%%%%%%%%%%%%%%%%%%%%%%%%%%%%%%%
\framesection{2.2 Détection d'objets}
\begin{rightframe}{DETR-ResNet}
\splitcols{%
  Object detection, a relatively old task in computer vision, is the task of identifying one or more pre-specified categories of objects inside the frame of an image. This combines what was originally two tasks: (1) classification and (2) localization, with most modern models for fixed classes finding that doing both together at the same time is best.

  We will use the DETR-RESNET-50 model, which was trained on the COCO 2017 dataset of 1000 categories tagged across over 100k images.\footcite{DBLP:journals/corr/abs-2005-12872} Like most object detection models, it produces confidence scores and rectangular bounding boxes that that object in question is predicted to be within. 
}{%
  \imgcaption{media/screenshot01.png}{Structure of the DETR-ResNet model from the original paper.}
}
\end{rightframe}


%%%%%%%%%%%%%%%%%%%%%%%%%%%%%%%%%%%%%%%%%%%%%%%%%%%%%%%%%%%%%%%%%%%%%%%%%%%%%
\begin{outlineframe}{Plan}
  \rmdim{2.1 Préparation des données}
  \rmdim{2.2 Détection d'objets}
  \rmitem{2.3 Détection à vocabulaire ouvert}
  \rmdim{2.4 Grounding DINO + TrOCR}
  \rmdim{2.5 Segment Anything}
  \rmdim{2.6 SAM + Grounding DINO}
  \rmdim{2.7 Estimation de profondeur}
  \rmdim{2.8 Plongements d'images}
  \rmdim{2.9 Modèles contrastifs}
  \rmdim{2.10 Détection et reconnaissance de visages}
  \rmdim{2.11 Estimation de pose}
  \rmdim{2.12 Modèles vision-langage (VLM)}
  \rmdim{2.13 VLM avec sortie structurée}
  \rmdim{2.14 Conclusion}
\end{outlineframe}

%%%%%%%%%%%%%%%%%%%%%%%%%%%%%%%%%%%%%%%%%%%%%%%%%%%%%%%%%%%%%%%%%%%%%%%%%%%%%
\framesection{2.3 Détection à vocabulaire ouvert}
\begin{rightframe}{Grounding DINO}
\splitcols{%
  A more recent version of object detection combines the idea of detection with a transformer-based language model to find arbitrarily described objects based on a textual description. We will use the Grounding DINO model, available with open-weights and runable with no specific architecture or hardware.\footcite{liu2023grounding}

  While the model is capable of running with any textual description, and produces reasonable results in high-resourced languages with common objects, it's best performance comes from describing the English names of the objects.
}{%
  \imgcaption{media/screenshot02.png}{Structure of the Grounding DINO model from the original paper.}
}
\end{rightframe}

%%%%%%%%%%%%%%%%%%%%%%%%%%%%%%%%%%%%%%%%%%%%%%%%%%%%%%%%%%%%%%%%%%%%%%%%%%%%%
\begin{outlineframe}{Plan}
  \rmdim{2.1 Préparation des données}
  \rmdim{2.2 Détection d'objets}
  \rmdim{2.3 Détection à vocabulaire ouvert}
  \rmitem{2.4 Grounding DINO + TrOCR}
  \rmdim{2.5 Segment Anything}
  \rmdim{2.6 SAM + Grounding DINO}
  \rmdim{2.7 Estimation de profondeur}
  \rmdim{2.8 Plongements d'images}
  \rmdim{2.9 Modèles contrastifs}
  \rmdim{2.10 Détection et reconnaissance de visages}
  \rmdim{2.11 Estimation de pose}
  \rmdim{2.12 Modèles vision-langage (VLM)}
  \rmdim{2.13 VLM avec sortie structurée}
  \rmdim{2.14 Conclusion}
\end{outlineframe}

%%%%%%%%%%%%%%%%%%%%%%%%%%%%%%%%%%%%%%%%%%%%%%%%%%%%%%%%%%%%%%%%%%%%%%%%%%%%%
\framesection{2.4 Grounding DINO + TrOCR}
\begin{rightframe}{Grounding DINO + TrOCR}
\splitcols{%
  Movie posters, like a like of historical images that we may be interested in studying, contain a lot of text, which is crutial to understanding the message of the poster. We can combine Grounding DINO with an OCR model to locate and then decode text found on each poster. We will use the TrOCR model, which is accessible through the same transformer's API that we have used with the other models so far.\footcite{li2022trocrtransformerbasedopticalcharacter}
}{%
  \imgcaption{media/screenshot03.png}{Structure of the TrOCR model.}
}
\end{rightframe}


%%%%%%%%%%%%%%%%%%%%%%%%%%%%%%%%%%%%%%%%%%%%%%%%%%%%%%%%%%%%%%%%%%%%%%%%%%%%%
\begin{outlineframe}{Plan}
  \rmdim{2.1 Préparation des données}
  \rmdim{2.2 Détection d'objets}
  \rmdim{2.3 Détection à vocabulaire ouvert}
  \rmdim{2.4 Grounding DINO + TrOCR}
  \rmitem{2.5 Segment Anything}
  \rmdim{2.6 SAM + Grounding DINO}
  \rmdim{2.7 Estimation de profondeur}
  \rmdim{2.8 Plongements d'images}
  \rmdim{2.9 Modèles contrastifs}
  \rmdim{2.10 Détection et reconnaissance de visages}
  \rmdim{2.11 Estimation de pose}
  \rmdim{2.12 Modèles vision-langage (VLM)}
  \rmdim{2.13 VLM avec sortie structurée}
  \rmdim{2.14 Conclusion}
\end{outlineframe}

%%%%%%%%%%%%%%%%%%%%%%%%%%%%%%%%%%%%%%%%%%%%%%%%%%%%%%%%%%%%%%%%%%%%%%%%%%%%%
\framesection{2.5 Segment Anything}
\begin{rightframe}{SAM: Segmentation à partir d'un point}
\splitcols{%
  The Segment Anything (SAM ; segmentation à partir d'un point) and its successors has had a lot publicity in the machine learning community.\footcite{kirillov2023segment} SAM allows for estimating the region, that is the specific set of pixels, that correspond to an object without needing to know beforehand what the object is. Previous attempts at segmentation assumed that we needed to know the object name first, but if we think about the way that we view things this makes some sense. We can identify the part of an image corresponding to a new animal even if we don't know what it is called.

  It is a very powful model, though as I think we will see, it's somewhat misunderstood because people are most familiar with the form in the image to the right from the original paper.
}{%
  \imgcaption{media/screenshot04.png}{Example of SAM in action.}
}
\end{rightframe}


%%%%%%%%%%%%%%%%%%%%%%%%%%%%%%%%%%%%%%%%%%%%%%%%%%%%%%%%%%%%%%%%%%%%%%%%%%%%%
\begin{outlineframe}{Plan}
  \rmdim{2.1 Préparation des données}
  \rmdim{2.2 Détection d'objets}
  \rmdim{2.3 Détection à vocabulaire ouvert}
  \rmdim{2.4 Grounding DINO + TrOCR}
  \rmdim{2.5 Segment Anything}
  \rmitem{2.6 SAM + Grounding DINO}
  \rmdim{2.7 Estimation de profondeur}
  \rmdim{2.8 Plongements d'images}
  \rmdim{2.9 Modèles contrastifs}
  \rmdim{2.10 Détection et reconnaissance de visages}
  \rmdim{2.11 Estimation de pose}
  \rmdim{2.12 Modèles vision-langage (VLM)}
  \rmdim{2.13 VLM avec sortie structurée}
  \rmdim{2.14 Conclusion}
\end{outlineframe}

%%%%%%%%%%%%%%%%%%%%%%%%%%%%%%%%%%%%%%%%%%%%%%%%%%%%%%%%%%%%%%%%%%%%%%%%%%%%%
\framesection{2.6 SAM + Grounding DINO}
\begin{rightframe}{Segmentation guidée par texte}
\splitcols{%
  As with the text example, we can combine a grounding model with SAM to generate the regions corresponding to any textual description of an object. This approach was written up and described as Grounded SAM, but we can implement it ourselves directly with the code we have already used.\footcite{ren2024grounded}
}{%
  \imgcaption{media/screenshot05.png}{Example of Grounded SAM.}
}
\end{rightframe}



%%%%%%%%%%%%%%%%%%%%%%%%%%%%%%%%%%%%%%%%%%%%%%%%%%%%%%%%%%%%%%%%%%%%%%%%%%%%%
\begin{outlineframe}{Plan}
  \rmdim{2.1 Préparation des données}
  \rmdim{2.2 Détection d'objets}
  \rmdim{2.3 Détection à vocabulaire ouvert}
  \rmdim{2.4 Grounding DINO + TrOCR}
  \rmdim{2.5 Segment Anything}
  \rmdim{2.6 SAM + Grounding DINO}
  \rmitem{2.7 Estimation de profondeur}
  \rmdim{2.8 Plongements d'images}
  \rmdim{2.9 Modèles contrastifs}
  \rmdim{2.10 Détection et reconnaissance de visages}
  \rmdim{2.11 Estimation de pose}
  \rmdim{2.12 Modèles vision-langage (VLM)}
  \rmdim{2.13 VLM avec sortie structurée}
  \rmdim{2.14 Conclusion}
\end{outlineframe}

%%%%%%%%%%%%%%%%%%%%%%%%%%%%%%%%%%%%%%%%%%%%%%%%%%%%%%%%%%%%%%%%%%%%%%%%%%%%%
\framesection{2.7 Estimation de profondeur}
\begin{rightframe}{Depth Anything}
\splitcols{%
  Our next model, called Depth Anything, is an attempt to measure the distance between any object and the camera.\footcite{yang2024depth} Some variants have the ability to provide exact metric measurements, though these are not very useful for our movie poster data, which does not consist of a single photograph from one angle but rather several elements arranged dynamicaly within a frame. So, we will stick with the faster original Depth Anything model here. 

  While exact distance to the camera does not have a specific definition for us here, we will see that it allows us to detangle the background and foreground of many of the movie posters.
}{%
  \imgcaption{media/screenshot06.png}{Architecture of the Depth Anything model, as described in the original paper.}
}
\end{rightframe}


%%%%%%%%%%%%%%%%%%%%%%%%%%%%%%%%%%%%%%%%%%%%%%%%%%%%%%%%%%%%%%%%%%%%%%%%%%%%%
\begin{outlineframe}{Plan}
  \rmdim{2.1 Préparation des données}
  \rmdim{2.2 Détection d'objets}
  \rmdim{2.3 Détection à vocabulaire ouvert}
  \rmdim{2.4 Grounding DINO + TrOCR}
  \rmdim{2.5 Segment Anything}
  \rmdim{2.6 SAM + Grounding DINO}
  \rmdim{2.7 Estimation de profondeur}
  \rmitem{2.8 Plongements d'images}
  \rmdim{2.9 Modèles contrastifs}
  \rmdim{2.10 Détection et reconnaissance de visages}
  \rmdim{2.11 Estimation de pose}
  \rmdim{2.12 Modèles vision-langage (VLM)}
  \rmdim{2.13 VLM avec sortie structurée}
  \rmdim{2.14 Conclusion}
\end{outlineframe}

%%%%%%%%%%%%%%%%%%%%%%%%%%%%%%%%%%%%%%%%%%%%%%%%%%%%%%%%%%%%%%%%%%%%%%%%%%%%%
\framesection{2.8 Plongements d'images}
\begin{rightframe}{DINOv2}
\splitcols{%
  Toutes les annotations vues jusqu'ici sont *interprétables* : on peut dire « il y a deux personnes sur cette affiche » ou « 30 \% de l'image est au premier plan ». Mais on peut aussi vouloir une représentation plus abstraite, qui résume l'image dans son ensemble en un vecteur de nombres — un *plongement* (embedding).

  Embeddings have been used extensively over the last decade and a half for working with both text and images. Well before the Grounding models that we have seen so far, these were the only way of working with categories that were defined ahead of time and specifically trained with a large dataset.

  Today we will use the DINOv2 model, the backbone for the Grounding DINO model we saw earlier.\footcite{oquab2023dinov2}
}{%
  \imgcaption{media/screenshot07.png}{Architecture of the DINOv2 model as described in the original paper.}
}
\end{rightframe}


%%%%%%%%%%%%%%%%%%%%%%%%%%%%%%%%%%%%%%%%%%%%%%%%%%%%%%%%%%%%%%%%%%%%%%%%%%%%%
\begin{outlineframe}{Plan}
  \rmdim{2.1 Préparation des données}
  \rmdim{2.2 Détection d'objets}
  \rmdim{2.3 Détection à vocabulaire ouvert}
  \rmdim{2.4 Grounding DINO + TrOCR}
  \rmdim{2.5 Segment Anything}
  \rmdim{2.6 SAM + Grounding DINO}
  \rmdim{2.7 Estimation de profondeur}
  \rmdim{2.8 Plongements d'images}
  \rmitem{2.9 Modèles contrastifs}
  \rmdim{2.10 Détection et reconnaissance de visages}
  \rmdim{2.11 Estimation de pose}
  \rmdim{2.12 Modèles vision-langage (VLM)}
  \rmdim{2.13 VLM avec sortie structurée}
  \rmdim{2.14 Conclusion}
\end{outlineframe}

%%%%%%%%%%%%%%%%%%%%%%%%%%%%%%%%%%%%%%%%%%%%%%%%%%%%%%%%%%%%%%%%%%%%%%%%%%%%%
\framesection{2.9 Modèles contrastifs}
\begin{rightframe}{SigLIP2}
\splitcols{%
  DINOv2 produit des plongements à partir d'images seules. Une autre
  famille de modèles, dits *contrastifs*, apprend conjointement à
  représenter images et textes dans un même espace. Le modèle le plus
  connu est CLIP ; SigLIP et sa version améliorée SigLIP2 en sont des
  évolutions plus performantes.\footcite{tschannen2025siglip}

  Today we will use SigLIP2 model, which allows us to embed images and text into the same high-dimensional space.
}{%
  \imgcaption{media/CLIP.png}{Conceptual view of the original CLIP.}
}
\end{rightframe}


%%%%%%%%%%%%%%%%%%%%%%%%%%%%%%%%%%%%%%%%%%%%%%%%%%%%%%%%%%%%%%%%%%%%%%%%%%%%%
\begin{outlineframe}{Plan}
  \rmdim{2.1 Préparation des données}
  \rmdim{2.2 Détection d'objets}
  \rmdim{2.3 Détection à vocabulaire ouvert}
  \rmdim{2.4 Grounding DINO + TrOCR}
  \rmdim{2.5 Segment Anything}
  \rmdim{2.6 SAM + Grounding DINO}
  \rmdim{2.7 Estimation de profondeur}
  \rmdim{2.8 Plongements d'images}
  \rmdim{2.9 Modèles contrastifs}
  \rmitem{2.10 Détection et reconnaissance de visages}
  \rmdim{2.11 Estimation de pose}
  \rmdim{2.12 Modèles vision-langage (VLM)}
  \rmdim{2.13 VLM avec sortie structurée}
  \rmdim{2.14 Conclusion}
\end{outlineframe}

%%%%%%%%%%%%%%%%%%%%%%%%%%%%%%%%%%%%%%%%%%%%%%%%%%%%%%%%%%%%%%%%%%%%%%%%%%%%%
\framesection{2.10 Détection et reconnaissance de visages}
\begin{rightframe}{SigLIP2}
\splitcols{%
  Les visages sont au cœur de la grammaire visuelle des affiches de
  films. Nous allons ici aller un cran plus loin avec la bibliothèque
  InsightFace, qui détecte les visages, estime l'âge et le genre
  apparents, et produit un plongement permettant de reconnaître si
  deux visages sont ceux de la même personne.\footcite{deng2021masked}

  Quelques mots de précaution s'imposent. L'âge et le genre prédits
  par ces modèles sont des estimations probabilistes, fondées sur les
  données d'entraînement utilisées. Ces données présentent des biais
  connus (sous-représentation de certaines populations, étiquettes
  binaires pour le genre, etc.), et il est important d'en tenir compte
  dans toute analyse.
}{%
  \imgcaption{media/facerec.png}{Visualization of face detection and recognition from the InsightFace website.}
}
\end{rightframe}


%%%%%%%%%%%%%%%%%%%%%%%%%%%%%%%%%%%%%%%%%%%%%%%%%%%%%%%%%%%%%%%%%%%%%%%%%%%%%
\begin{outlineframe}{Plan}
  \rmdim{2.1 Préparation des données}
  \rmdim{2.2 Détection d'objets}
  \rmdim{2.3 Détection à vocabulaire ouvert}
  \rmdim{2.4 Grounding DINO + TrOCR}
  \rmdim{2.5 Segment Anything}
  \rmdim{2.6 SAM + Grounding DINO}
  \rmdim{2.7 Estimation de profondeur}
  \rmdim{2.8 Plongements d'images}
  \rmdim{2.9 Modèles contrastifs}
  \rmdim{2.10 Détection et reconnaissance de visages}
  \rmitem{2.11 Estimation de pose}
  \rmdim{2.12 Modèles vision-langage (VLM)}
  \rmdim{2.13 VLM avec sortie structurée}
  \rmdim{2.14 Conclusion}
\end{outlineframe}

%%%%%%%%%%%%%%%%%%%%%%%%%%%%%%%%%%%%%%%%%%%%%%%%%%%%%%%%%%%%%%%%%%%%%%%%%%%%%
\framesection{2.11 Estimation de pose}
\begin{rightframe}{ViTPose}
\splitcols{%
  Au-delà de la simple détection d'une personne, on peut chercher à
  caractériser sa *pose* : position de la tête, des bras, des jambes.
  L'estimation de pose consiste à localiser un ensemble de points-clés
  anatomiques (nez, yeux, épaules, coudes, etc.) sur le corps humain.

  Notre pipeline procède en deux étapes : on détecte d'abord les
  personnes avec DETR, puis on applique ViTPose, un modèle dédié à
  l'estimation de pose, sur chaque personne détectée.\footcite{xu2022vitpose}
}{%
  \imgcaption{media/screenshot08.png}{Example of ViTPose in action from the original paper.}
}
\end{rightframe}


%%%%%%%%%%%%%%%%%%%%%%%%%%%%%%%%%%%%%%%%%%%%%%%%%%%%%%%%%%%%%%%%%%%%%%%%%%%%%
\begin{outlineframe}{Plan}
  \rmdim{2.1 Préparation des données}
  \rmdim{2.2 Détection d'objets}
  \rmdim{2.3 Détection à vocabulaire ouvert}
  \rmdim{2.4 Grounding DINO + TrOCR}
  \rmdim{2.5 Segment Anything}
  \rmdim{2.6 SAM + Grounding DINO}
  \rmdim{2.7 Estimation de profondeur}
  \rmdim{2.8 Plongements d'images}
  \rmdim{2.9 Modèles contrastifs}
  \rmdim{2.10 Détection et reconnaissance de visages}
  \rmdim{2.11 Estimation de pose}
  \rmitem{2.12 Modèles vision-langage (VLM)}
  \rmdim{2.13 VLM avec sortie structurée}
  \rmdim{2.14 Conclusion}
\end{outlineframe}

%%%%%%%%%%%%%%%%%%%%%%%%%%%%%%%%%%%%%%%%%%%%%%%%%%%%%%%%%%%%%%%%%%%%%%%%%%%%%
\framesection{2.12 Modèles vision-langage (VLM)}
\begin{rightframe}{ViTPose}
\splitcols{%
  Les modèles vus jusqu'ici produisent des sorties structurées : des
  boîtes, des masques, des vecteurs. Les modèles vision-langage (VLM)
  font quelque chose de différent : on leur fournit une image et une
  question en langage naturel, et ils répondent en langage naturel.

  We could use many open-weight models, either locally or through a service such as OpenRouter. For simplicity, today we will use the the OpenAI API. This is nice because it can be adapted to call local models through tools such as Ollama or LMStudio.
}{%
  \imgcaption{media/qwen-omni.png}{The Qwen2.5-Omni model, one of the most popular open-weight general-purpose VLM models.}
}
\end{rightframe}

%%%%%%%%%%%%%%%%%%%%%%%%%%%%%%%%%%%%%%%%%%%%%%%%%%%%%%%%%%%%%%%%%%%%%%%%%%%%%
\begin{outlineframe}{Plan}
  \rmdim{2.1 Préparation des données}
  \rmdim{2.2 Détection d'objets}
  \rmdim{2.3 Détection à vocabulaire ouvert}
  \rmdim{2.4 Grounding DINO + TrOCR}
  \rmdim{2.5 Segment Anything}
  \rmdim{2.6 SAM + Grounding DINO}
  \rmdim{2.7 Estimation de profondeur}
  \rmdim{2.8 Plongements d'images}
  \rmdim{2.9 Modèles contrastifs}
  \rmdim{2.10 Détection et reconnaissance de visages}
  \rmdim{2.11 Estimation de pose}
  \rmdim{2.12 Modèles vision-langage (VLM)}
  \rmitem{2.13 VLM avec sortie structurée}
  \rmdim{2.14 Conclusion}
\end{outlineframe}

%%%%%%%%%%%%%%%%%%%%%%%%%%%%%%%%%%%%%%%%%%%%%%%%%%%%%%%%%%%%%%%%%%%%%%%%%%%%%
\framesection{2.13 VLM avec sortie structurée}
\begin{rightframe}{Pydantic}
\splitcols{%
  We can further make use of VLM models by forcing them to produce structured data. The most popular way to do this, at least as of 2026, is through Pydantic. This is a tool that describes the schema that we want in the output data and then verifies that the VLM has correctly structured the data as requested.
}{%
  \imgcaption{media/screenshot08.png}{Homepage for the Pydantic validation tool.}
}
\end{rightframe}


%%%%%%%%%%%%%%%%%%%%%%%%%%%%%%%%%%%%%%%%%%%%%%%%%%%%%%%%%%%%%%%%%%%%%%%%%%%%%
\begin{outlineframe}{Plan}
  \rmdim{2.1 Préparation des données}
  \rmdim{2.2 Détection d'objets}
  \rmdim{2.3 Détection à vocabulaire ouvert}
  \rmdim{2.4 Grounding DINO + TrOCR}
  \rmdim{2.5 Segment Anything}
  \rmdim{2.6 SAM + Grounding DINO}
  \rmdim{2.7 Estimation de profondeur}
  \rmdim{2.8 Plongements d'images}
  \rmdim{2.9 Modèles contrastifs}
  \rmdim{2.10 Détection et reconnaissance de visages}
  \rmdim{2.11 Estimation de pose}
  \rmdim{2.12 Modèles vision-langage (VLM)}
  \rmdim{2.13 VLM avec sortie structurée}
  \rmitem{2.14 Conclusion}
\end{outlineframe}

%%%%%%%%%%%%%%%%%%%%%%%%%%%%%%%%%%%%%%%%%%%%%%%%%%%%%%%%%%%%%%%%%%%%%%%%%%%%%
\framesection{2.14 Conclusion}
\begin{rightframe}{Next Steps}
We have seen a variety of different models in this tutorial. These touch on many of the types of cutting-edge models available in mid-2026 for working with still, visual image data.

Both the specific models and, at a slower pace the tasks, will continue to change and update over the coming years. But as can be seen here, new models build on previous ones and often do not entirely replace the older models.

It's important to keep up with developments in this space going forward, but the tools here should be useful for years to come.
\end{rightframe}

%%%%%%%%%%%%%%%%%%%%%%%%%%%%%%%%%%%%%%%%%%%%%%%%%%%%%%%%%%%%%%%%%%%%%%%%%%%%%
\framesection{}
\begin{frame}[allowframebreaks]{References}
  \printbibliography[heading=none]
\end{frame}

\end{document}